\chapter{Maximum-Likelihood, MAP und Bayessche Schätzung}

\section{Ziel dieses Kapitels}

Im vorherigen Kapitel haben wir Schätzer (estimators) allgemein betrachtet.
Wir haben gesehen, dass ein Schätzer eine Funktion der Daten ist und dass man Schätzer nach Bias, Varianz, mittlerem quadratischem Fehler und Konsistenz beurteilen kann.

In diesem Kapitel betrachten wir drei zentrale Prinzipien zur Parameterschätzung:

\begin{itemize}
    \item Maximum-Likelihood-Schätzung (maximum likelihood estimation),
    \item Maximum-a-posteriori-Schätzung (maximum a-posteriori estimation),
    \item Bayessche Schätzung (Bayesian estimation).
\end{itemize}

\begin{conceptbox}
Maximum-Likelihood-Schätzung (maximum likelihood estimation) verwendet nur die Daten.

Maximum-a-posteriori-Schätzung (maximum a-posteriori estimation) verwendet Daten und einen Prior.

Bayessche Schätzung (Bayesian estimation) berechnet nicht nur einen besten Parameterwert, sondern eine ganze Posterior-Verteilung über Parameter.
\end{conceptbox}

Nach diesem Kapitel solltest du folgende Fragen beantworten können:

\begin{itemize}
    \item Was ist eine Likelihood?
    \item Warum ist die Likelihood keine Wahrscheinlichkeitsverteilung über Parameter?
    \item Wie funktioniert Maximum-Likelihood-Schätzung?
    \item Warum verwendet man die Log-Likelihood?
    \item Wie hängt Maximum-Likelihood mit Verlustfunktionen zusammen?
    \item Warum führt Gaußsches Rauschen zu quadratischem Fehler?
    \item Was ist ein Prior?
    \item Was ist ein Posterior?
    \item Wie wird MAP aus Bayes' Regel hergeleitet?
    \item Warum kann MAP als regularisierte Maximum-Likelihood-Schätzung verstanden werden?
    \item Was unterscheidet MLE, MAP und Bayesian Estimation?
    \item Was ist eine prädiktive Verteilung (predictive distribution)?
\end{itemize}

\section{Grundsituation der Parameterschätzung}

Wir betrachten ein Modell mit Parameter \(\theta\).
Der Parameter \(\theta\) kann eine einzelne Zahl, ein Vektor oder auch eine Menge mehrerer Parameter sein.

Beispiele:

\[
    \theta = \mu
\]

für den Mittelwert einer Normalverteilung,

\[
    \theta = (\mu,\sigma^2)
\]

für Mittelwert und Varianz einer Normalverteilung,

oder

\[
    \theta = \vec{w}
\]

für die Gewichte eines linearen Modells.

Wir beobachten Daten

\[
    \mathcal{D}
    =
    \{x_1,x_2,\dots,x_N\}.
\]

Oft nehmen wir an, dass die Daten unabhängig und identisch verteilt sind (independent and identically distributed), kurz i.i.d.:

\[
    x_1,\dots,x_N
    \overset{\mathrm{i.i.d.}}{\sim}
    p(x \mid \theta).
\]

Das bedeutet:

\[
    p(\mathcal{D} \mid \theta)
    =
    p(x_1,\dots,x_N \mid \theta)
    =
    \prod_{n=1}^{N}
    p(x_n \mid \theta).
\]

\begin{definitionbox}
Die Größe

\[
    p(\mathcal{D} \mid \theta)
\]

beschreibt, wie wahrscheinlich oder wie dicht die beobachteten Daten \(\mathcal{D}\) unter dem Parameter \(\theta\) sind.
\end{definitionbox}

\section{Likelihood}

Die Likelihood ist dieselbe mathematische Größe wie \(p(\mathcal{D} \mid \theta)\), aber mit anderer Interpretation.

\begin{definitionbox}
Für feste beobachtete Daten \(\mathcal{D}\) ist die Likelihood (likelihood) eine Funktion des Parameters \(\theta\):

\[
    L(\theta)
    =
    p(\mathcal{D} \mid \theta).
\]

Sie misst, wie plausibel verschiedene Parameterwerte \(\theta\) die beobachteten Daten machen.
\end{definitionbox}

Wichtig ist die Perspektive:

\[
    p(\mathcal{D} \mid \theta)
\]

als Wahrscheinlichkeitsmodell:
\(\theta\) ist fest, \(\mathcal{D}\) ist zufällig.

\[
    L(\theta)
\]

als Likelihood:
\(\mathcal{D}\) ist fest beobachtet, \(\theta\) wird variiert.

\begin{notebox}
Die Likelihood \(L(\theta)\) ist im Allgemeinen keine Wahrscheinlichkeitsverteilung über \(\theta\).

Das heißt:
Es gilt im Allgemeinen nicht

\[
    \int L(\theta)\,d\theta = 1.
\]

Oder im diskreten Fall:

\[
    \sum_{\theta} L(\theta) = 1.
\]

Die Likelihood ist nur bis auf ihre relative Größe für verschiedene Parameterwerte wichtig.
\end{notebox}

\section{Maximum-Likelihood-Schätzung}

Die Maximum-Likelihood-Schätzung (maximum likelihood estimation, MLE) wählt den Parameterwert, der die beobachteten Daten unter dem Modell maximal plausibel macht.

\begin{definitionbox}
Der Maximum-Likelihood-Schätzer ist definiert als

\[
    \hat{\theta}_{\mathrm{ML}}
    =
    \arg\max_{\theta}
    p(\mathcal{D} \mid \theta).
\]

Mit der Likelihood-Schreibweise:

\[
    \hat{\theta}_{\mathrm{ML}}
    =
    \arg\max_{\theta}
    L(\theta).
\]
\end{definitionbox}

Für i.i.d. Daten gilt:

\[
    p(\mathcal{D} \mid \theta)
    =
    \prod_{n=1}^{N}
    p(x_n \mid \theta).
\]

Daher:

\[
    \hat{\theta}_{\mathrm{ML}}
    =
    \arg\max_{\theta}
    \prod_{n=1}^{N}
    p(x_n \mid \theta).
\]

\begin{conceptbox}
Maximum-Likelihood bedeutet:

\[
    \text{Wähle } \theta \text{ so, dass die beobachteten Daten möglichst wahrscheinlich sind.}
\]
\end{conceptbox}

\section{Log-Likelihood}

Produkte vieler kleiner Wahrscheinlichkeiten oder Dichten sind numerisch und algebraisch unpraktisch.
Deshalb maximiert man meistens den Logarithmus der Likelihood.

\begin{definitionbox}
Die Log-Likelihood (log-likelihood) ist

\[
    \ell(\theta)
    =
    \log L(\theta)
    =
    \log p(\mathcal{D} \mid \theta).
\]
\end{definitionbox}

Da der Logarithmus streng monoton wachsend ist, gilt:

\[
    \arg\max_{\theta}
    L(\theta)
    =
    \arg\max_{\theta}
    \log L(\theta).
\]

Für i.i.d. Daten:

\[
    \ell(\theta)
    =
    \log
    \prod_{n=1}^{N}
    p(x_n \mid \theta).
\]

Mit der Logarithmusregel

\[
    \log \prod_{n=1}^{N} a_n
    =
    \sum_{n=1}^{N} \log a_n
\]

folgt:

\[
    \ell(\theta)
    =
    \sum_{n=1}^{N}
    \log p(x_n \mid \theta).
\]

\begin{conceptbox}
Die Log-Likelihood wandelt ein Produkt in eine Summe um:

\[
    \prod_{n=1}^{N}
    p(x_n \mid \theta)
    \quad
    \longrightarrow
    \quad
    \sum_{n=1}^{N}
    \log p(x_n \mid \theta).
\]

Summen sind leichter abzuleiten, numerisch stabiler und passen direkt zu Verlustfunktionen.
\end{conceptbox}

\section{Negative Log-Likelihood}

In maschinellem Lernen formuliert man viele Probleme als Minimierung einer Verlustfunktion (loss function).
Da Maximum-Likelihood ein Maximierungsproblem ist, verwendet man oft die negative Log-Likelihood.

\begin{definitionbox}
Die negative Log-Likelihood (negative log-likelihood, NLL) ist

\[
    \mathcal{L}(\theta)
    =
    -\ell(\theta)
    =
    -\log p(\mathcal{D} \mid \theta).
\]
\end{definitionbox}

Maximum-Likelihood ist äquivalent zur Minimierung der negativen Log-Likelihood:

\[
    \hat{\theta}_{\mathrm{ML}}
    =
    \arg\max_{\theta}
    \log p(\mathcal{D} \mid \theta)
\]

ist äquivalent zu

\[
    \hat{\theta}_{\mathrm{ML}}
    =
    \arg\min_{\theta}
    \left[
        -\log p(\mathcal{D} \mid \theta)
    \right].
\]

Für i.i.d. Daten:

\[
    -\log p(\mathcal{D} \mid \theta)
    =
    -\sum_{n=1}^{N}
    \log p(x_n \mid \theta).
\]

Also:

\[
    \mathcal{L}(\theta)
    =
    \sum_{n=1}^{N}
    -\log p(x_n \mid \theta).
\]

\begin{conceptbox}
Viele Verlustfunktionen im maschinellen Lernen sind negative Log-Likelihoods.

Das bedeutet:
Man kann viele Trainingsziele probabilistisch interpretieren.
\end{conceptbox}

\section{Standardrezept für MLE-Herleitungen}

Viele MLE-Aufgaben folgen demselben Muster.

\begin{examtipbox}
Standardvorgehen für Maximum-Likelihood-Schätzung:

\begin{enumerate}
    \item Modell \(p(x \mid \theta)\) aufschreiben.
    \item Likelihood für alle Daten bilden:
    \[
        L(\theta)
        =
        \prod_{n=1}^{N}
        p(x_n \mid \theta).
    \]
    \item Log-Likelihood bilden:
    \[
        \ell(\theta)
        =
        \sum_{n=1}^{N}
        \log p(x_n \mid \theta).
    \]
    \item Nach \(\theta\) ableiten.
    \item Ableitung gleich \(0\) setzen.
    \item Nach \(\theta\) auflösen.
    \item Prüfen, ob ein Maximum vorliegt und ob die Lösung zulässig ist.
\end{enumerate}
\end{examtipbox}

\section{Beispiel: Bernoulli-Verteilung}

Die Bernoulli-Verteilung (Bernoulli distribution) modelliert eine binäre Zufallsvariable:

\[
    X \in \{0,1\}.
\]

Der Parameter ist

\[
    \pi \in [0,1].
\]

Die Verteilung lautet:

\[
    p(x \mid \pi)
    =
    \pi^x(1-\pi)^{1-x}.
\]

Dabei gilt:

\[
    p(X=1 \mid \pi) = \pi,
\]

und

\[
    p(X=0 \mid \pi) = 1-\pi.
\]

Wir beobachten Daten

\[
    x_1,\dots,x_N
    \in
    \{0,1\}.
\]

\subsection{Likelihood}

Für i.i.d. Daten ist die Likelihood:

\[
    L(\pi)
    =
    \prod_{n=1}^{N}
    \pi^{x_n}
    (1-\pi)^{1-x_n}.
\]

\subsection{Log-Likelihood}

Die Log-Likelihood ist:

\[
    \ell(\pi)
    =
    \sum_{n=1}^{N}
    \log
    \left[
        \pi^{x_n}
        (1-\pi)^{1-x_n}
    \right].
\]

Mit Logarithmusregeln:

\[
    \ell(\pi)
    =
    \sum_{n=1}^{N}
    \left[
        x_n \log \pi
        +
        (1-x_n)\log(1-\pi)
    \right].
\]

\subsection{Ableitung}

Ableiten nach \(\pi\):

\[
    \frac{\partial \ell}{\partial \pi}
    =
    \sum_{n=1}^{N}
    \left[
        \frac{x_n}{\pi}
        -
        \frac{1-x_n}{1-\pi}
    \right].
\]

Setze

\[
    S
    =
    \sum_{n=1}^{N}
    x_n.
\]

Dann gilt:

\[
    \sum_{n=1}^{N}
    (1-x_n)
    =
    N-S.
\]

Damit:

\[
    \frac{\partial \ell}{\partial \pi}
    =
    \frac{S}{\pi}
    -
    \frac{N-S}{1-\pi}.
\]

Setze die Ableitung gleich \(0\):

\[
    \frac{S}{\pi}
    -
    \frac{N-S}{1-\pi}
    =
    0.
\]

Also:

\[
    \frac{S}{\pi}
    =
    \frac{N-S}{1-\pi}.
\]

Kreuzmultiplizieren:

\[
    S(1-\pi)
    =
    \pi(N-S).
\]

Ausmultiplizieren:

\[
    S - S\pi
    =
    \pi N - \pi S.
\]

Damit:

\[
    S = \pi N.
\]

Also:

\[
    \hat{\pi}_{\mathrm{ML}}
    =
    \frac{S}{N}
    =
    \frac{1}{N}
    \sum_{n=1}^{N}
    x_n.
\]

\begin{conceptbox}
Der Maximum-Likelihood-Schätzer für den Bernoulli-Parameter ist der relative Anteil der Einsen:

\[
    \hat{\pi}_{\mathrm{ML}}
    =
    \frac{1}{N}
    \sum_{n=1}^{N}
    x_n.
\]
\end{conceptbox}

\section{Bernoulli-MLE als Klassifikationsverlust}

Die negative Log-Likelihood der Bernoulli-Verteilung ist:

\[
    -\ell(\pi)
    =
    -\sum_{n=1}^{N}
    \left[
        x_n \log \pi
        +
        (1-x_n)\log(1-\pi)
    \right].
\]

Für ein Modell, das für jeden Datenpunkt eine Wahrscheinlichkeit

\[
    y_n = p(t_n=1 \mid \vec{x}_n)
\]

ausgibt, erhält man den Verlust

\[
    \mathcal{L}
    =
    -\sum_{n=1}^{N}
    \left[
        t_n \log y_n
        +
        (1-t_n)\log(1-y_n)
    \right].
\]

\begin{definitionbox}
Der Ausdruck

\[
    -\left[
        t \log y
        +
        (1-t)\log(1-y)
    \right]
\]

heißt binäre Kreuzentropie (binary cross-entropy).

Sie ist die negative Log-Likelihood eines Bernoulli-Modells.
\end{definitionbox}

\begin{notebox}
Die binäre Kreuzentropie ist also nicht nur eine willkürliche Verlustfunktion.
Sie entsteht direkt aus Maximum-Likelihood für Bernoulli-verteilte Zielwerte.
\end{notebox}

\section{Beispiel: Normalverteilung mit bekanntem \(\sigma^2\)}

Nun betrachten wir kontinuierliche Daten.
Seien

\[
    X_1,\dots,X_N
    \overset{\mathrm{i.i.d.}}{\sim}
    \mathcal{N}(\mu,\sigma^2),
\]

wobei \(\sigma^2\) bekannt ist und \(\mu\) geschätzt werden soll.

Die Dichte ist:

\[
    p(x_n \mid \mu,\sigma^2)
    =
    \frac{1}{\sqrt{2\pi\sigma^2}}
    \exp\left(
        -\frac{(x_n-\mu)^2}{2\sigma^2}
    \right).
\]

\subsection{Likelihood}

Die Likelihood ist:

\[
    L(\mu)
    =
    \prod_{n=1}^{N}
    \frac{1}{\sqrt{2\pi\sigma^2}}
    \exp\left(
        -\frac{(x_n-\mu)^2}{2\sigma^2}
    \right).
\]

\subsection{Log-Likelihood}

Die Log-Likelihood ist:

\[
    \ell(\mu)
    =
    \sum_{n=1}^{N}
    \left[
        -\frac{1}{2}\log(2\pi\sigma^2)
        -
        \frac{(x_n-\mu)^2}{2\sigma^2}
    \right].
\]

Also:

\[
    \ell(\mu)
    =
    -\frac{N}{2}\log(2\pi\sigma^2)
    -
    \frac{1}{2\sigma^2}
    \sum_{n=1}^{N}
    (x_n-\mu)^2.
\]

Der erste Term hängt nicht von \(\mu\) ab.
Daher maximieren wir \(\ell(\mu)\), indem wir minimieren:

\[
    \sum_{n=1}^{N}
    (x_n-\mu)^2.
\]

\begin{conceptbox}
Für Gaußsches Rauschen führt Maximum-Likelihood auf quadratische Fehler.

Das ist der Grund, warum Least Squares und Normalverteilung eng zusammenhängen.
\end{conceptbox}

\subsection{Ableitung}

Wir minimieren:

\[
    E(\mu)
    =
    \sum_{n=1}^{N}
    (x_n-\mu)^2.
\]

Ableiten:

\[
    \frac{\partial E}{\partial \mu}
    =
    \sum_{n=1}^{N}
    2(x_n-\mu)(-1).
\]

Also:

\[
    \frac{\partial E}{\partial \mu}
    =
    -2
    \sum_{n=1}^{N}
    (x_n-\mu).
\]

Setze gleich \(0\):

\[
    -2
    \sum_{n=1}^{N}
    (x_n-\mu)
    =
    0.
\]

Damit:

\[
    \sum_{n=1}^{N}
    (x_n-\mu)
    =
    0.
\]

Also:

\[
    \sum_{n=1}^{N}
    x_n
    -
    \sum_{n=1}^{N}
    \mu
    =
    0.
\]

Da \(\mu\) konstant ist:

\[
    \sum_{n=1}^{N}
    \mu
    =
    N\mu.
\]

Somit:

\[
    \sum_{n=1}^{N}
    x_n
    -
    N\mu
    =
    0.
\]

Daraus folgt:

\[
    \hat{\mu}_{\mathrm{ML}}
    =
    \frac{1}{N}
    \sum_{n=1}^{N}
    x_n.
\]

\begin{definitionbox}
Für normalverteilte Daten mit bekannter Varianz ist der Maximum-Likelihood-Schätzer des Mittelwerts das Stichprobenmittel:

\[
    \hat{\mu}_{\mathrm{ML}}
    =
    \frac{1}{N}
    \sum_{n=1}^{N}
    x_n.
\]
\end{definitionbox}

\section{Normalverteilung mit unbekanntem \(\mu\) und \(\sigma^2\)}

Jetzt seien beide Parameter unbekannt:

\[
    \theta = (\mu,\sigma^2).
\]

Die Daten seien

\[
    X_1,\dots,X_N
    \overset{\mathrm{i.i.d.}}{\sim}
    \mathcal{N}(\mu,\sigma^2).
\]

Die Log-Likelihood ist:

\[
    \ell(\mu,\sigma^2)
    =
    -\frac{N}{2}\log(2\pi\sigma^2)
    -
    \frac{1}{2\sigma^2}
    \sum_{n=1}^{N}
    (x_n-\mu)^2.
\]

Maximum-Likelihood liefert:

\[
    \hat{\mu}_{\mathrm{ML}}
    =
    \frac{1}{N}
    \sum_{n=1}^{N}
    x_n,
\]

und

\[
    \hat{\sigma}_{\mathrm{ML}}^2
    =
    \frac{1}{N}
    \sum_{n=1}^{N}
    (x_n-\hat{\mu}_{\mathrm{ML}})^2.
\]

\begin{notebox}
Der MLE für die Varianz der Normalverteilung verwendet den Nenner \(N\), nicht \(N-1\).

Er ist daher bei endlichem \(N\) verzerrt:

\[
    \mathbb{E}[\hat{\sigma}_{\mathrm{ML}}^2]
    =
    \frac{N-1}{N}
    \sigma^2.
\]
\end{notebox}

\section{Regression als Maximum-Likelihood}

Maximum-Likelihood erklärt auch die lineare Regression (linear regression).

Wir betrachten Daten

\[
    (\vec{x}_1,t_1),\dots,(\vec{x}_N,t_N).
\]

Ein lineares Modell ist:

\[
    y(\vec{x},\vec{w})
    =
    \vec{w}^{\top}\vec{x}.
\]

Wir nehmen an, dass die Zielwerte mit Gaußschem Rauschen entstehen:

\[
    t_n
    =
    \vec{w}^{\top}\vec{x}_n
    +
    \epsilon_n,
\]

wobei

\[
    \epsilon_n
    \sim
    \mathcal{N}(0,\sigma^2).
\]

Dann gilt:

\[
    p(t_n \mid \vec{x}_n,\vec{w},\sigma^2)
    =
    \mathcal{N}
    \left(
        t_n
        \mid
        \vec{w}^{\top}\vec{x}_n,
        \sigma^2
    \right).
\]

Für i.i.d. Daten ist:

\[
    p(\vec{t} \mid X,\vec{w},\sigma^2)
    =
    \prod_{n=1}^{N}
    \mathcal{N}
    \left(
        t_n
        \mid
        \vec{w}^{\top}\vec{x}_n,
        \sigma^2
    \right).
\]

Die negative Log-Likelihood ist bis auf Konstanten:

\[
    \mathcal{L}(\vec{w})
    =
    \frac{1}{2\sigma^2}
    \sum_{n=1}^{N}
    \left(
        t_n
        -
        \vec{w}^{\top}\vec{x}_n
    \right)^2.
\]

Da \(\frac{1}{2\sigma^2}\) positiv und konstant ist, ist das äquivalent zu:

\[
    E(\vec{w})
    =
    \frac{1}{2}
    \sum_{n=1}^{N}
    \left(
        t_n
        -
        \vec{w}^{\top}\vec{x}_n
    \right)^2.
\]

\begin{conceptbox}
Lineare Regression mit quadratischem Fehler ist Maximum-Likelihood unter der Annahme von Gaußschem Rauschen.

\[
    \text{Gaußsches Rauschen}
    \quad
    \Longrightarrow
    \quad
    \text{quadratischer Fehler}.
\]
\end{conceptbox}

\section{Prior}

Maximum-Likelihood verwendet nur die Daten.
In vielen Situationen haben wir aber Vorwissen über Parameter.
Dieses Vorwissen wird durch einen Prior beschrieben.

\begin{definitionbox}
Ein Prior \(p(\theta)\) ist eine Wahrscheinlichkeitsverteilung über Parameter vor Betrachtung der aktuellen Daten.
\end{definitionbox}

Beispiele:

\[
    p(\mu)
    =
    \mathcal{N}(\mu \mid 0,\tau^2)
\]

bedeutet:
Vor den Daten halten wir Werte von \(\mu\) nahe \(0\) für plausibler als sehr große Werte.

Für einen Gewichtsvektor \(\vec{w}\) kann ein Prior sein:

\[
    p(\vec{w})
    =
    \mathcal{N}
    \left(
        \vec{w}
        \mid
        \vec{0},
        \alpha^{-1}I
    \right).
\]

Dieser Prior bevorzugt kleine Gewichte.

\begin{notebox}
Der Prior ist kein Trick, sondern eine Modellannahme.

Er beschreibt, welche Parameterwerte vor Beobachtung der Daten plausibel sind.
\end{notebox}

\section{Posterior}

Wenn Daten beobachtet werden, aktualisiert Bayes' Regel die Verteilung über Parameter.

\[
    p(\theta \mid \mathcal{D})
    =
    \frac{
        p(\mathcal{D} \mid \theta)p(\theta)
    }{
        p(\mathcal{D})
    }.
\]

\begin{definitionbox}
Der Posterior \(p(\theta \mid \mathcal{D})\) ist die Verteilung über Parameter nach Beobachtung der Daten.
\end{definitionbox}

Die einzelnen Bestandteile sind:

\[
    \underbrace{p(\theta \mid \mathcal{D})}_{\text{Posterior}}
    =
    \frac{
        \underbrace{p(\mathcal{D} \mid \theta)}_{\text{Likelihood}}
        \underbrace{p(\theta)}_{\text{Prior}}
    }{
        \underbrace{p(\mathcal{D})}_{\text{Evidence}}
    }.
\]

Die Evidence ist:

\[
    p(\mathcal{D})
    =
    \int
    p(\mathcal{D} \mid \theta)
    p(\theta)
    \,d\theta.
\]

Im diskreten Fall:

\[
    p(\mathcal{D})
    =
    \sum_{\theta}
    p(\mathcal{D} \mid \theta)
    p(\theta).
\]

\begin{conceptbox}
Bayes' Regel für Parameter lautet:

\[
    \text{Posterior}
    =
    \frac{
        \text{Likelihood}
        \cdot
        \text{Prior}
    }{
        \text{Evidence}
    }.
\]

Oder proportional:

\[
    p(\theta \mid \mathcal{D})
    \propto
    p(\mathcal{D} \mid \theta)p(\theta).
\]
\end{conceptbox}

\section{Maximum-a-posteriori-Schätzung}

Die Maximum-a-posteriori-Schätzung (maximum a-posteriori estimation, MAP) wählt den Parameterwert, der den Posterior maximiert.

\begin{definitionbox}
Der MAP-Schätzer ist

\[
    \hat{\theta}_{\mathrm{MAP}}
    =
    \arg\max_{\theta}
    p(\theta \mid \mathcal{D}).
\]
\end{definitionbox}

Mit Bayes' Regel:

\[
    p(\theta \mid \mathcal{D})
    =
    \frac{
        p(\mathcal{D} \mid \theta)p(\theta)
    }{
        p(\mathcal{D})
    }.
\]

Da \(p(\mathcal{D})\) nicht von \(\theta\) abhängt, ist der Nenner für die Maximierung irrelevant:

\[
    \hat{\theta}_{\mathrm{MAP}}
    =
    \arg\max_{\theta}
    p(\mathcal{D} \mid \theta)p(\theta).
\]

Mit Logarithmus:

\[
    \hat{\theta}_{\mathrm{MAP}}
    =
    \arg\max_{\theta}
    \left[
        \log p(\mathcal{D} \mid \theta)
        +
        \log p(\theta)
    \right].
\]

\begin{conceptbox}
MAP maximiert die Summe aus Log-Likelihood und Log-Prior:

\[
    \log p(\theta \mid \mathcal{D})
    =
    \log p(\mathcal{D} \mid \theta)
    +
    \log p(\theta)
    -
    \log p(\mathcal{D}).
\]

Für die Maximierung nach \(\theta\) kann \(\log p(\mathcal{D})\) ignoriert werden.
\end{conceptbox}

\section{MLE als Spezialfall von MAP}

Wenn der Prior konstant ist, dann ist MAP dasselbe wie MLE.

Angenommen:

\[
    p(\theta) = \text{konstant}.
\]

Dann:

\[
    \hat{\theta}_{\mathrm{MAP}}
    =
    \arg\max_{\theta}
    p(\mathcal{D} \mid \theta)p(\theta).
\]

Da \(p(\theta)\) konstant ist:

\[
    \hat{\theta}_{\mathrm{MAP}}
    =
    \arg\max_{\theta}
    p(\mathcal{D} \mid \theta)
    =
    \hat{\theta}_{\mathrm{ML}}.
\]

\begin{notebox}
Man sagt oft:
MLE entspricht MAP mit einem flachen oder uninformativem Prior.

Das ist eine nützliche Intuition, aber mathematisch muss man bei kontinuierlichen Parametern vorsichtig sein, weil ein wirklich konstanter Prior auf ganz \(\mathbb{R}\) nicht normalisierbar ist.
\end{notebox}

\section{MAP und Regularisierung}

MAP ist besonders wichtig, weil es Regularisierung (regularization) probabilistisch erklärt.

Wir betrachten wieder ein lineares Modell:

\[
    y(\vec{x},\vec{w})
    =
    \vec{w}^{\top}\vec{x}.
\]

Die Likelihood unter Gaußschem Rauschen führt auf den quadratischen Fehler:

\[
    \frac{1}{2\sigma^2}
    \sum_{n=1}^{N}
    \left(
        t_n
        -
        \vec{w}^{\top}\vec{x}_n
    \right)^2.
\]

Nun nehmen wir einen Gauß-Prior auf \(\vec{w}\) an:

\[
    p(\vec{w})
    =
    \mathcal{N}
    \left(
        \vec{w}
        \mid
        \vec{0},
        \alpha^{-1}I
    \right).
\]

Die Dichte ist proportional zu:

\[
    p(\vec{w})
    \propto
    \exp\left(
        -\frac{\alpha}{2}
        \vec{w}^{\top}\vec{w}
    \right).
\]

Der negative Log-Prior ist daher bis auf Konstanten:

\[
    -\log p(\vec{w})
    =
    \frac{\alpha}{2}
    \vec{w}^{\top}\vec{w}
    +
    \text{const}.
\]

Die negative Log-Posterior ist also bis auf Konstanten:

\[
    \mathcal{L}_{\mathrm{MAP}}(\vec{w})
    =
    \frac{1}{2\sigma^2}
    \sum_{n=1}^{N}
    \left(
        t_n
        -
        \vec{w}^{\top}\vec{x}_n
    \right)^2
    +
    \frac{\alpha}{2}
    \vec{w}^{\top}\vec{w}.
\]

Multipliziert man mit der positiven Konstanten \(\sigma^2\), erhält man ein äquivalentes Optimierungsproblem:

\[
    E(\vec{w})
    =
    \frac{1}{2}
    \sum_{n=1}^{N}
    \left(
        t_n
        -
        \vec{w}^{\top}\vec{x}_n
    \right)^2
    +
    \frac{\lambda}{2}
    \vec{w}^{\top}\vec{w},
\]

mit

\[
    \lambda = \alpha\sigma^2.
\]

\begin{conceptbox}
Ein Gauß-Prior auf Gewichte führt zu \(L_2\)-Regularisierung:

\[
    \frac{\lambda}{2}
    \vec{w}^{\top}\vec{w}.
\]

Das ist die probabilistische Interpretation von Ridge Regression.
\end{conceptbox}

\section{Laplace-Prior und \(L_1\)-Regularisierung}

Nicht nur Gauß-Priors sind möglich.
Ein Laplace-Prior auf Gewichte hat die Form:

\[
    p(w_i)
    \propto
    \exp(-\lambda |w_i|).
\]

Für einen Gewichtsvektor mit unabhängigen Komponenten:

\[
    p(\vec{w})
    \propto
    \exp\left(
        -\lambda
        \sum_{i}
        |w_i|
    \right).
\]

Dann ist der negative Log-Prior proportional zu:

\[
    \lambda
    \sum_i
    |w_i|.
\]

\begin{conceptbox}
Ein Laplace-Prior führt zu \(L_1\)-Regularisierung:

\[
    \lambda
    \sum_i
    |w_i|.
\]

\(L_1\)-Regularisierung fördert sparse Lösungen, also Lösungen, bei denen viele Gewichte exakt oder nahezu \(0\) sind.
\end{conceptbox}

\section{Bayessche Schätzung}

MLE und MAP liefern jeweils einen einzelnen Parameterwert.
Bayessche Schätzung (Bayesian estimation) betrachtet dagegen die ganze Posterior-Verteilung.

\begin{definitionbox}
Bei Bayesscher Schätzung wird der Parameter \(\theta\) selbst als Zufallsvariable behandelt.

Nach Beobachtung der Daten beschreibt der Posterior

\[
    p(\theta \mid \mathcal{D})
\]

die Unsicherheit über \(\theta\).
\end{definitionbox}

Das Ziel ist also nicht nur:

\[
    \hat{\theta}
\]

zu finden, sondern die gesamte Verteilung

\[
    p(\theta \mid \mathcal{D})
\]

zu berechnen oder zu approximieren.

\begin{conceptbox}
Bayesian Estimation behält Unsicherheit über Parameter bei.

Das ist besonders wichtig, wenn wenige Daten vorhanden sind oder wenn Entscheidungen die Unsicherheit berücksichtigen sollen.
\end{conceptbox}

\section{Punktschätzer aus dem Posterior}

Auch in der Bayesschen Statistik kann man aus dem Posterior einen einzelnen Punktschätzer ableiten.

\subsection{Posterior-Mittelwert}

Der Posterior-Mittelwert ist:

\[
    \hat{\theta}_{\mathrm{Bayes}}
    =
    \mathbb{E}[\theta \mid \mathcal{D}].
\]

Im kontinuierlichen Fall:

\[
    \hat{\theta}_{\mathrm{Bayes}}
    =
    \int
    \theta
    p(\theta \mid \mathcal{D})
    \,d\theta.
\]

\subsection{Posterior-Modus}

Der Posterior-Modus ist der Wert mit größter Posterior-Dichte:

\[
    \hat{\theta}_{\mathrm{mode}}
    =
    \arg\max_{\theta}
    p(\theta \mid \mathcal{D}).
\]

Das ist genau der MAP-Schätzer:

\[
    \hat{\theta}_{\mathrm{mode}}
    =
    \hat{\theta}_{\mathrm{MAP}}.
\]

\begin{notebox}
MAP ist ein Punktschätzer aus dem Posterior.

Bayessche Schätzung als Ganzes ist aber mehr als MAP:
Sie betrachtet die gesamte Posterior-Verteilung.
\end{notebox}

\section{Prädiktive Verteilung}

In maschinellem Lernen wollen wir oft nicht nur Parameter schätzen, sondern Vorhersagen für neue Daten machen.

Sei \(x_*\) ein neuer Eingabewert und \(t_*\) der unbekannte Zielwert.
Wenn \(\theta\) unbekannt ist, sollte man die Unsicherheit über \(\theta\) in der Vorhersage berücksichtigen.

\begin{definitionbox}
Die prädiktive Verteilung (predictive distribution) ist

\[
    p(t_* \mid x_*,\mathcal{D})
    =
    \int
    p(t_* \mid x_*,\theta)
    p(\theta \mid \mathcal{D})
    \,d\theta.
\]
\end{definitionbox}

Diese Formel mittelt Vorhersagen über alle möglichen Parameterwerte, gewichtet mit deren Posterior-Wahrscheinlichkeit.

\begin{conceptbox}
Bayessche Vorhersage integriert über Parameterunsicherheit:

\[
    \text{Vorhersage}
    =
    \text{Mittel über alle plausiblen Parameter}.
\]

Das unterscheidet Bayesian Prediction von einer reinen Punktschätzung.
\end{conceptbox}

\section{Beispiel: Beta-Bernoulli-Modell}

Ein klassisches Beispiel für Bayessche Schätzung ist das Beta-Bernoulli-Modell.

Wir beobachten binäre Daten:

\[
    x_n \in \{0,1\}.
\]

Die Likelihood ist Bernoulli:

\[
    p(x_n \mid \pi)
    =
    \pi^{x_n}
    (1-\pi)^{1-x_n}.
\]

Für den Parameter \(\pi\) verwenden wir einen Beta-Prior:

\[
    p(\pi)
    =
    \mathrm{Beta}(\pi \mid a,b).
\]

Die Beta-Verteilung hat die Form:

\[
    \mathrm{Beta}(\pi \mid a,b)
    =
    \frac{\Gamma(a+b)}{\Gamma(a)\Gamma(b)}
    \pi^{a-1}
    (1-\pi)^{b-1},
\]

für

\[
    0 \leq \pi \leq 1.
\]

\subsection{Likelihood}

Für \(N\) Datenpunkte:

\[
    p(\mathcal{D} \mid \pi)
    =
    \prod_{n=1}^{N}
    \pi^{x_n}
    (1-\pi)^{1-x_n}.
\]

Setze

\[
    S
    =
    \sum_{n=1}^{N}
    x_n.
\]

Dann ist

\[
    \sum_{n=1}^{N}
    (1-x_n)
    =
    N-S.
\]

Also:

\[
    p(\mathcal{D} \mid \pi)
    =
    \pi^S
    (1-\pi)^{N-S}.
\]

\subsection{Posterior}

Nach Bayes:

\[
    p(\pi \mid \mathcal{D})
    \propto
    p(\mathcal{D} \mid \pi)p(\pi).
\]

Einsetzen:

\[
    p(\pi \mid \mathcal{D})
    \propto
    \pi^S
    (1-\pi)^{N-S}
    \pi^{a-1}
    (1-\pi)^{b-1}.
\]

Zusammenfassen der Exponenten:

\[
    p(\pi \mid \mathcal{D})
    \propto
    \pi^{a+S-1}
    (1-\pi)^{b+N-S-1}.
\]

Das ist wieder eine Beta-Verteilung:

\[
    p(\pi \mid \mathcal{D})
    =
    \mathrm{Beta}(\pi \mid a+S,b+N-S).
\]

\begin{conceptbox}
Beim Beta-Bernoulli-Modell ist der Posterior wieder eine Beta-Verteilung:

\[
    \mathrm{Beta}(a,b)
    \quad
    \longrightarrow
    \quad
    \mathrm{Beta}(a+S,b+N-S).
\]

Der Parameter \(a\) zählt sozusagen vorherige Erfolge, \(b\) vorherige Misserfolge.
\end{conceptbox}

\section{Konjugierte Priors}

Das Beta-Bernoulli-Beispiel zeigt eine wichtige Idee.

\begin{definitionbox}
Ein Prior heißt konjugierter Prior (conjugate prior) für eine Likelihood, wenn der Posterior zur selben Verteilungsfamilie gehört wie der Prior.
\end{definitionbox}

Beim Beta-Bernoulli-Modell:

\[
    \text{Prior}
    =
    \mathrm{Beta},
\]

\[
    \text{Likelihood}
    =
    \mathrm{Bernoulli},
\]

\[
    \text{Posterior}
    =
    \mathrm{Beta}.
\]

\begin{notebox}
Konjugierte Priors sind mathematisch praktisch, weil der Posterior analytisch berechnet werden kann.

In komplexeren ML-Modellen ist der Posterior oft nicht analytisch berechenbar.
Dann braucht man Approximationen.
\end{notebox}

\section{Vergleich von MLE, MAP und Bayesscher Schätzung}

\[
\begin{array}{c|c|c|c}
\text{Methode}
&
\text{Ziel}
&
\text{Prior}
&
\text{Ergebnis}
\\
\hline
\text{MLE}
&
\max_{\theta} p(\mathcal{D}\mid\theta)
&
\text{nein}
&
\hat{\theta}_{\mathrm{ML}}
\\
\hline
\text{MAP}
&
\max_{\theta} p(\theta\mid\mathcal{D})
&
\text{ja}
&
\hat{\theta}_{\mathrm{MAP}}
\\
\hline
\text{Bayesian}
&
p(\theta\mid\mathcal{D})
&
\text{ja}
&
\text{Posterior-Verteilung}
\end{array}
\]

\begin{conceptbox}
MLE:
\[
    \hat{\theta}_{\mathrm{ML}}
    =
    \arg\max_{\theta}
    p(\mathcal{D}\mid\theta).
\]

MAP:
\[
    \hat{\theta}_{\mathrm{MAP}}
    =
    \arg\max_{\theta}
    p(\mathcal{D}\mid\theta)p(\theta).
\]

Bayesian:
\[
    p(\theta\mid\mathcal{D})
    =
    \frac{
        p(\mathcal{D}\mid\theta)p(\theta)
    }{
        p(\mathcal{D})
    }.
\]
\end{conceptbox}

\section{Interpretation bei wenigen und vielen Daten}

Der Unterschied zwischen MLE, MAP und Bayesian Estimation ist besonders bei kleinen Datensätzen wichtig.

\subsection{Wenige Daten}

Wenn nur wenige Daten vorhanden sind, kann MLE stark schwanken.
Ein Prior kann dann helfen, unplausible Parameterwerte zu vermeiden.

\[
    \text{wenige Daten}
    \quad
    \Rightarrow
    \quad
    \text{Prior hat starken Einfluss}.
\]

\subsection{Viele Daten}

Wenn sehr viele Daten vorhanden sind, dominiert oft die Likelihood.
Dann wird der Einfluss des Priors relativ kleiner.

\[
    \text{viele Daten}
    \quad
    \Rightarrow
    \quad
    \text{Likelihood dominiert}.
\]

\begin{notebox}
Das bedeutet nicht, dass der Prior bei vielen Daten immer egal ist.
Aber in vielen einfachen regulären Modellen wird sein Einfluss im Vergleich zur Datenmenge kleiner.
\end{notebox}

\section{Zusammenhang zu Overfitting}

Maximum-Likelihood kann bei flexiblen Modellen zu Überanpassung (overfitting) führen.
Das passiert, wenn Parameter so gewählt werden, dass sie die Trainingsdaten sehr gut erklären, aber auf neuen Daten schlecht generalisieren.

MAP kann dem entgegenwirken, wenn der Prior komplexe oder extreme Parameterwerte bestraft.

\begin{conceptbox}
Regularisierung kann probabilistisch als Prior interpretiert werden.

\[
    \text{Regularisierung}
    \quad
    \Longleftrightarrow
    \quad
    \text{Prior auf Parameter}.
\]
\end{conceptbox}

Beispiel:

\[
    \frac{\lambda}{2}
    \vec{w}^{\top}\vec{w}
\]

bestraft große Gewichte.
Das entspricht einem Gauß-Prior auf \(\vec{w}\).

\section{Häufige Fehler}

\begin{examtipbox}
Typische Fehler bei MLE, MAP und Bayesian Estimation:

\begin{enumerate}
    \item Likelihood \(p(\mathcal{D}\mid\theta)\) mit Posterior \(p(\theta\mid\mathcal{D})\) verwechseln.
    \item Denken, dass die Likelihood eine Wahrscheinlichkeitsverteilung über \(\theta\) sein muss.
    \item Bei MLE den Logarithmus falsch verwenden.
    \item Vergessen, dass Maximieren der Log-Likelihood äquivalent zum Maximieren der Likelihood ist.
    \item Negative Log-Likelihood und Log-Likelihood mit falschem Vorzeichen verwenden.
    \item Bei MAP den Prior vergessen.
    \item Die Evidence \(p(\mathcal{D})\) bei der Maximierung nach \(\theta\) unnötig mitoptimieren.
    \item MAP mit vollständiger Bayesscher Schätzung gleichsetzen.
    \item Quadratischen Fehler als rein willkürlichen Verlust betrachten, obwohl er aus Gaußschem Rauschen folgt.
    \item \(L_2\)-Regularisierung nicht als Gauß-Prior erkennen.
\end{enumerate}
\end{examtipbox}

\section{Mini-Aufgaben}

\subsection{Aufgabe 1: Bernoulli-MLE}

Gegeben sind die Daten

\[
    x_1=1,
    \quad
    x_2=0,
    \quad
    x_3=1,
    \quad
    x_4=1,
    \quad
    x_5=0.
\]

Berechne den Maximum-Likelihood-Schätzer für \(\pi\).

\begin{examplebox}
Für Bernoulli-Daten gilt:

\[
    \hat{\pi}_{\mathrm{ML}}
    =
    \frac{1}{N}
    \sum_{n=1}^{N}
    x_n.
\]

Hier ist:

\[
    N=5
\]

und

\[
    \sum_{n=1}^{5}
    x_n
    =
    1+0+1+1+0
    =
    3.
\]

Also:

\[
    \hat{\pi}_{\mathrm{ML}}
    =
    \frac{3}{5}
    =
    0{,}6.
\]
\end{examplebox}

\subsection{Aufgabe 2: Negative Log-Likelihood}

Für ein Bernoulli-Modell sei

\[
    t=1
\]

und das Modell sagt

\[
    y = 0{,}8.
\]

Berechne die negative Log-Likelihood für diesen einzelnen Datenpunkt.

\begin{examplebox}
Die negative Log-Likelihood ist:

\[
    -\left[
        t\log y
        +
        (1-t)\log(1-y)
    \right].
\]

Mit \(t=1\):

\[
    -\left[
        1\cdot\log(0{,}8)
        +
        0\cdot\log(0{,}2)
    \right]
    =
    -\log(0{,}8).
\]

Also:

\[
    \mathcal{L}
    =
    -\log(0{,}8).
\]
\end{examplebox}

\subsection{Aufgabe 3: MAP mit konstantem Prior}

Zeige, dass MAP mit konstantem Prior zu MLE wird.

\begin{examplebox}
MAP ist:

\[
    \hat{\theta}_{\mathrm{MAP}}
    =
    \arg\max_{\theta}
    p(\mathcal{D}\mid\theta)p(\theta).
\]

Wenn

\[
    p(\theta)=c
\]

konstant ist, dann:

\[
    \hat{\theta}_{\mathrm{MAP}}
    =
    \arg\max_{\theta}
    c\,p(\mathcal{D}\mid\theta).
\]

Da \(c\) positiv und konstant ist, ändert es die Stelle des Maximums nicht:

\[
    \hat{\theta}_{\mathrm{MAP}}
    =
    \arg\max_{\theta}
    p(\mathcal{D}\mid\theta)
    =
    \hat{\theta}_{\mathrm{ML}}.
\]
\end{examplebox}

\subsection{Aufgabe 4: Beta-Bernoulli-Posterior}

Ein Prior sei

\[
    \pi \sim \mathrm{Beta}(2,3).
\]

In den Daten gibt es

\[
    S=4
\]

Einsen und

\[
    N-S=6
\]

Nullen.
Berechne den Posterior.

\begin{examplebox}
Beim Beta-Bernoulli-Modell gilt:

\[
    \mathrm{Beta}(a,b)
    \longrightarrow
    \mathrm{Beta}(a+S,b+N-S).
\]

Hier ist:

\[
    a=2,
    \quad
    b=3,
    \quad
    S=4,
    \quad
    N-S=6.
\]

Also:

\[
    \pi \mid \mathcal{D}
    \sim
    \mathrm{Beta}(2+4,3+6).
\]

Damit:

\[
    \pi \mid \mathcal{D}
    \sim
    \mathrm{Beta}(6,9).
\]
\end{examplebox}

\section{Zusammenfassung}

\begin{itemize}
    \item Die Likelihood ist:
    \[
        L(\theta)
        =
        p(\mathcal{D}\mid\theta).
    \]
    \item Bei i.i.d. Daten faktorisiert sie:
    \[
        p(\mathcal{D}\mid\theta)
        =
        \prod_{n=1}^{N}
        p(x_n\mid\theta).
    \]
    \item Die Log-Likelihood ist:
    \[
        \ell(\theta)
        =
        \sum_{n=1}^{N}
        \log p(x_n\mid\theta).
    \]
    \item MLE maximiert die Likelihood:
    \[
        \hat{\theta}_{\mathrm{ML}}
        =
        \arg\max_{\theta}
        p(\mathcal{D}\mid\theta).
    \]
    \item Negative Log-Likelihood ist eine Verlustfunktion:
    \[
        \mathcal{L}(\theta)
        =
        -\log p(\mathcal{D}\mid\theta).
    \]
    \item Bernoulli-MLE ergibt den Anteil der Einsen:
    \[
        \hat{\pi}_{\mathrm{ML}}
        =
        \frac{1}{N}
        \sum_{n=1}^{N}
        x_n.
    \]
    \item Gaußsche Likelihood führt zu quadratischem Fehler.
    \item Ein Prior \(p(\theta)\) beschreibt Vorwissen über Parameter.
    \item Der Posterior ist:
    \[
        p(\theta\mid\mathcal{D})
        =
        \frac{
            p(\mathcal{D}\mid\theta)p(\theta)
        }{
            p(\mathcal{D})
        }.
    \]
    \item MAP maximiert den Posterior:
    \[
        \hat{\theta}_{\mathrm{MAP}}
        =
        \arg\max_{\theta}
        p(\theta\mid\mathcal{D}).
    \]
    \item MAP ist äquivalent zu:
    \[
        \hat{\theta}_{\mathrm{MAP}}
        =
        \arg\max_{\theta}
        p(\mathcal{D}\mid\theta)p(\theta).
    \]
    \item Ein Gauß-Prior auf Gewichte führt zu \(L_2\)-Regularisierung.
    \item Ein Laplace-Prior führt zu \(L_1\)-Regularisierung.
    \item Bayessche Schätzung betrachtet die gesamte Posterior-Verteilung \(p(\theta\mid\mathcal{D})\).
    \item Die prädiktive Verteilung integriert über Parameterunsicherheit:
    \[
        p(t_* \mid x_*,\mathcal{D})
        =
        \int
        p(t_* \mid x_*,\theta)
        p(\theta\mid\mathcal{D})
        \,d\theta.
    \]
\end{itemize}

\section{Typische Prüfungsfragen}

\begin{examtipbox}
Mögliche Prüfungsfragen zu diesem Kapitel:

\begin{enumerate}
    \item Was ist eine Likelihood?
    \item Warum ist die Likelihood keine Wahrscheinlichkeitsverteilung über Parameter?
    \item Definiere den Maximum-Likelihood-Schätzer.
    \item Warum verwendet man die Log-Likelihood?
    \item Was ist die negative Log-Likelihood?
    \item Leite den MLE für den Parameter einer Bernoulli-Verteilung her.
    \item Warum führt Bernoulli-Likelihood zur binären Kreuzentropie?
    \item Leite den MLE für den Mittelwert einer Normalverteilung mit bekannter Varianz her.
    \item Warum führt Gaußsches Rauschen zu quadratischem Fehler?
    \item Was ist ein Prior?
    \item Was ist ein Posterior?
    \item Leite MAP aus Bayes' Regel her.
    \item Warum kann die Evidence bei MAP ignoriert werden?
    \item Wann ist MAP gleich MLE?
    \item Wie hängt MAP mit Regularisierung zusammen?
    \item Warum führt ein Gauß-Prior zu \(L_2\)-Regularisierung?
    \item Warum führt ein Laplace-Prior zu \(L_1\)-Regularisierung?
    \item Was ist der Unterschied zwischen MAP und vollständiger Bayesscher Schätzung?
    \item Was ist eine prädiktive Verteilung?
    \item Erkläre das Beta-Bernoulli-Modell.
\end{enumerate}
\end{examtipbox}